% !TeX root = ../papers/nc-p3.tex
\section{三項中心の第二三次曲線と三次スクロールの因子}
\label{sec:ternary-center-cubic-scroll}
\begin{theorem}\label{thm:ternary-center-cubic-scroll}
曲面射影と二直線列のない標準正則窓で、中心 $p$ の第二曲線がねじれ三次曲線なら、
$w(-,x_1(s),x_2(t),-)=[s,t]A(t)$ の二項二次行列は一般階数3である。
反対階数10では、この場合は中央支持3の正則三行三列鉛筆を持つ順中心に限られ、
第三曲線は直線である。第二曲線上の順側欠損は二次列二本と一次列一本の五次行列式で与えられる。
特に特異三項鉛筆の中心の第二曲線は二次である。逆順でも成立する。
\end{theorem}
\begin{proof}[左一次関係のある四空間]
$U=\ker p$ と置く。二項中心の排除により第一直線の一般中央支持は4、
逆支持は3である。従って列の空間 $D=A(t)V_3^*\subset U\otimes S_2$ は四次元で $U$ の全体を使う。
一般階数1なら二項多項式環の一意分解で多項式列と行の積となるが、三次元と四次元を
張る両因子には少なくとも次数2と3が必要である。そこで一般階数2と仮定する。

$R_0$ を第一直線に制限する写像は単射ではない。単射なら一般の $x_2(t)$ の後続に
$\ker A(t)$ の全直線が入る。前節までと同じく制限像は $[-,t]\otimes A(t)z$ の包であり、
一般階数2、完全な $U$ 支持から、その消去式は原始な一次行一つである。
よって制限写像の核は一次元。適合基底を取ると全列は
\[
          A(t)z=uQ_z(t)+v(t)L_z(t),\qquad
          v(t)=t_0v_0+t_1v_1
\]
となり、$u,v_0,v_1$ は独立、$Q_z$ は二次、$L_z$ は一次である。
$D$ は $S_2\oplus S_1$ 内の超平面である。$L$ の像が一次元なら
$D=S_2u+k n(t)v(t)$ で先のII型となるので排除される
（定理\ref{thm:rank-ten-cubic-kernel-divisor}の因子の証明は両標準形の下で反対階数を使わない）。
像が二次元なら超平面の式は $\lambda(Q)+\mu(L)=0$、$\lambda\ne0$ である。
$\lambda$ を $S_1$ 上の対称二形式 $\lambda(lm)$ と見る。
階数1で $\mu$ がこの形式の像に入る場合、$v(t)$ に一次式倍の $u$ を加えて $\mu$ を消せる。
$\ker\lambda$ の二次式は全て固定一次式の倍であり、先のI型となってやはり排除される。
\end{proof}
\begin{proof}[入射曲面上の辺と次の小行列式]
残りは $\lambda$ の階数2、または階数1で $\mu$ がその像に入らない場合である。
$A(t)z=0$ の解を二項変数 $y$ を用いて
\[
      Q_z(y)=[t,y]M(y),\quad L_z(y)=[t,y]b,\qquad M=M_0 y_0+M_1y_1
\]
と書く。$[M_0:M_1:b]\in\mathbb P^2$ と $t\in\mathbb P^1$ の間の
$\lambda([t,y]M)+b\mu([t,y])=0$ は階数2の双一次式である。
実際第一項の $M$ での係数行の写像は $\lambda$ と同階数で、階数1の場合には $b$ の項が独立な一次式を与える。
この曲面 $T$ は一点 $o$ での $\mathbb P^2$ の吹上げである。
二つの独立一次式を消す $o$ の外でその比が $t$ を定め、入射式がこのグラフの閉包を与える。
$F,H$ をそれぞれ $\mathbb P^1,\mathbb P^2$ の超平面類の引戻し、
$J$ を例外曲線とすれば $J=H-F$, $J^2=-1$, $F\cdot J=1$ である。
$z$ の座標は基点のない $H+F$ の切断であり、$b\ne0$ では $L_z$ の根が $t$ を復元するため
その像は曲面である。

前節の $E(t)z$ は同じく位置2と3の辺を与える固定核直線に値を取り、
$T$ 上では類 $H+4F$ の有効因子 $\Delta$ を切る。恒等零は可動中心の禁止に反する。
各ファイバーで一次なので、その水平部分は第三点を与える重複度1の切断である。
各素因子の一般点でその重複度まで次点 $x_4$ が取れることは、前節で使った逆支持3と
直交性による局所商環の議論の通りである。

$p$ からの射影の座標 $(u,v_0,v_1)$ において $h=(b,-M_0,-M_1)$ と置き、位置0への線型持上げも $h$ と書く。
$A(y)z$ との対合は全て零なので $M_j u(x_4)+b v_j(x_4)=0$ ($j=0,1$)。
従って $b$ と $(u(x_4),v_0(x_4),v_1(x_4))\wedge h$ の積は零である。
$R_3$ の $A(y)$ の三係数は独立である。実際 $v_0,v_1$ の方向での係数には
二つの独立な $L_z$ と全ての一次式の積が現れる。残る三式 $B_j$ について
$B_j(z,p),B_j(z,h)$ を二列とする。小行列の類は $3H+2F$ である。
$\Delta$ の素因子の重複度を $n$、$b$ のそこでの位数を $e$ とすれば、その小行列は
少なくとも $n-e$ の位数を持つ。実際離散付値環の商で上記の外積は一様化変数の
$n-e$ 乗を法に消える。$h$ の一成分は単元なので、この商で $x_4$ は $p,h$ の包に入る。
従って二列の間に単元成分を持つ関係がある。
\end{proof}
\begin{proof}[例外曲線での余分な直線の排除]
階数2の $\lambda$ では $b=0$ の直線は $o$ を通らないので、その引戻しは既約な水平曲線である。
この曲線が $\Delta$ の水平切断になることはない。その上で $z$ の像は
$L_z=0$, $Q_z\in\ker\lambda$ の射影直線で、パラメータ $t$ を根に持つ唯一の二次式が $Q_z$ となる。
この直線への射は次数2である。二式の鉛筆 $\ker\lambda$ が共通根を持てば
$\lambda$ はその点の評価の倍で階数1になってしまうからである。
これは全ての動く曲線のパラメータが像へ双有理であることに反する。

従ってこの場合は小行列から $\Delta$ 全体を除けて、残る三成分は類 $2H-2F=2J$ の切断である。
$H^0(\mathcal O(kJ))$ は $k\ge0$ で一次元である（$J^2=-1$ により順に $J$ を除く）。
よって全小行列が零か一定ベクトル倍になる。前者なら一般の $z$ で $p,h$ の包に
後続が存在して曲面射影が生じる。後者なら第一列 $(B_j(z,p))_j$ が固定平面に入る。
しかし $p$ の逆支持は3で $A(y)$ は $p$ を消し、$T$ の $z$ 像は全 $\mathbb P^3$ を張る。
最後の事実は任意の一般の $[t,y](M,b)$ の包、あるいは次の直接計算で分かる。
線型式が $z$ の像で消えればその双一次引戻しは入射式の倍である。
$S_2^*\oplus S_1^*$ から $(t,M,b)$ の双一次式へのこの代入は単射なので
元の式は $D$ の定義式の倍である。従って後者も不可能である。

階数1の残りの場合は $b=0$ が $o$ を通り、その引戻しは被約和 $J+F_e$ となる。
$F_e$ は $\lambda([t,y]M)$ が全 $M$ で零になる唯一のファイバーである。
$J$ と $F_e$ の $z$ 像は同じ直線 $Z=\{L_z=0\}$ であり、どちらの制限も同型である。
実際 $F_e$ 上では $[t,y]$ を固定して $M$ が動き、$J$ 上では $M$ が固定され $[t,y]$ が動く。
$\lambda$ が一つの点での評価なので、どちらもその固定一次因子を持つ二次式全体になる。
$F_e$ が $\Delta$ の成分でなければ、$\Delta$ から除く必要があるのは高々 $J$ 一つで、
小行列の残余類は $2J$ または $3J$。既に排除した定数比例の矛盾となる。
$F_e$ が成分なら、$\Delta$ から $F_e$ を一回と、必要なら $J$ を一回差し引いた因子を除ける。
残余の類は $kJ+F$、$k=2,3$ のいずれかである。$J$ への次数 $1-k<0$ のため
全残余成分が $J$ 上で消える。従って $J$ 上の一般点で $z$ は $p$ と固定した $h(o)$ の結ぶ直線に後続を持つ。
そこでは $b=0$ だが、$p,h(o)$ はいずれも $A(y)z$ を消すのでこの結論に影響しない。
他方、$F_e$ が辺の因子なので $Z$ は順中心 $x_2(e)$ の後続直線そのものである。
一般点の後続は一意なのでこの中心の第二曲線が直線となり、二直線列の禁止に反する。
これで一般階数2を全て排除した。
\end{proof}
\begin{proof}[反対階数十での五次欠損]
中央支持が4なら前定理が第二三次を排除するので、現在は3である。
$\dim S_1D\le10-3=7$ より、四次元 $D$ はある非零一次列 $L$ の
$t_0 L,t_1 L$ を含む（一次関係を割る前節までの計算）。
残る二列を $B,C$ とすれば $A(t)=(t_0L,t_1L,B,C)$。
一般階数3のため核直線は $(-t_1,t_0,0,0)$ であり、第三曲線は直線となる。
この場合 $[-,t]\otimes A(t)z$ は六次元を張るので $R_0$ の第一直線への制限は同型である。
よって $x_2$ の順側欠損は $A(t)$ の三次小行列の共通因子、すなわち非零五次式 $\det(L,B,C)$ が与える。
特異三項鉛筆なら命題\ref{prop:singular-ternary-center-projections}の下界2に反する。
\end{proof}

\begin{corollary}\label{cor:rank-ten-uniform-corridor-types}
反対階数10の曲面射影のない標準正則窓で分岐があれば、全ての列の第二曲線が二次であるか、
全てが次数列 $(1,3,1)$ で両中心とも正則三項鉛筆を持つかのいずれかである。
後者では各位相に順中心・逆中心はそれぞれ高々一つであり、その存在位相は両向きで一致し、
二位置移動で保たれる。前者では順中心も逆中心も全四位相に存在する。
\end{corollary}
\begin{proof}
任意の標準窓の分岐列で、直線からねじれ三次曲線へ写る辺を考える。
その直線の二項座標と次位置の四座標との双一次式で辺上消えるものは、
ねじれ三次の三つの標準一次関係の包に入る。実際直線と三次のパラメータは
同じ正規化で、三次は完全系を使う。この辺の関係の直線への制限は一般点の順評価が
階数3なのでちょうどこの三空間である。直線の一般点 $x$ の縮約を
$\sum K_j(x)\otimes B_j$ と書くと、三つの $K_j(x)$ は独立である。
その逆支持は $B_j$ の固定支持の和に等しく、従って一意な一般前点は一定である。
よってこの直線は列の最初であり、逆にねじれ三次に直線が続けば最後の直線である。
上の定理より第二三次に続くのは直線なので列は $(1,3,1)$ で、逆中心にも同じ定理を用いる。

第二二次の場合は最初の直線の唯一の順中心の鉛筆を読む。その二項行列 $A(t)$ と
この中心の鉛筆の一致は定理\ref{thm:rank-ten-conic-after-center}の最初の標準二列の証明による。
中央支持3でこの三行三列が正則なら $\ker A(t)$ は一般に一定の第四方向だけであるから
第三点が一定になり矛盾する。従ってこの中心は特異三項または四次元であり、既にどちらも
第二曲線が二次と分かっている。これを反復すれば全位相に第二二次の順中心がある。
一つの列で最後から第二曲線が三次なら逆向きの結果からその列自身の次数列が $(1,3,1)$ となる。従って第二二次だけの場合、逆向きでも同じ反復ができる。

一方、$(1,3,1)$ の順中心 $p$ を位置0に取ると同位相に他の順中心はない。
実際その後続直線は $L_p$ と交わらない。$p$ の正則三項鉛筆をその直線上で評価すれば
二項パラメータは全 $\mathbb P^1$ を動くので、一般に三行列は可逆であり、
第二点は $p$ の中央支持の消去点に固定されてしまう。
よって二次の中心鎖とは共存できない。三次の各列の両中心は同位相にあり、
第二三次上での順側欠損は5なので反対位相にも順中心がある。逆向きも同じである。
\end{proof}
